\documentclass[11pt]{article}
\usepackage[T1]{fontenc}
\usepackage{lmodern}
\usepackage[a4paper,margin=1in]{geometry}
\usepackage{amsmath}
\usepackage{booktabs}
\usepackage{graphicx}
\usepackage{xcolor}
\usepackage{caption}
\usepackage[hidelinks]{hyperref}

\newcommand{\figbox}[2]{%
  \IfFileExists{#1}{\includegraphics[width=#2\linewidth]{#1}}%
  {\fbox{\parbox[c][4.2cm][c]{#2\linewidth}{\centering\small\color{gray}%
   [\texttt{#1}]\\[2pt] generated by \texttt{studyA\_gradients.py};\\
   drop the PNG in \texttt{StudyA-figs/} and recompile.}}}}

\title{\textbf{Study A: Spatial Preservation vs.\ Parameter Count}\\[2pt]
\large New trials vs.\ old, model depths, and why skip connections help}
\author{Abdulrahman Mahmoud --- DLoc lightweight localization}
\date{\today}

\begin{document}
\maketitle

\noindent This report answers the four things requested: (1)~the \textbf{new
trials vs.\ the old models}, isolating the accuracy gained \emph{because we kept
the spatial relation between adjacent pixels}; (2)~a \textbf{snapshot of the
gradient descent} during training; (3)~\textbf{why skip connections raise
accuracy}, including the gradient-flow reasoning; and (4)~the \textbf{depth of
every model} used. All Study~A numbers use the \emph{mean} localization error
(the primary metric; the median is pixel-quantized to a $0.1\sqrt{n}$ lattice
and cannot resolve sub-decimetre differences).

%============================================================
\section{New trials vs.\ old --- the gain from spatial preservation}

Study~A builds one configurable student and toggles a single knob at a time, so
any two rows are a controlled comparison. The encoder is held at the
ImageNet-pretrained MobileNetV2; the capacity axis is the \emph{decoder width},
so extra parameters are added \emph{without} touching the pretraining.

\begin{table}[h]
\centering
\caption{Study~A results (mean / median / P90 localization error, metres).
``no-skip'' = plain MobileNetV2 decoder (no spatial preservation); ``skip'' =
U-Net skip connections; ``ASPP'' = multi-scale bottleneck.}
\setlength{\tabcolsep}{5pt}\small
\begin{tabular}{llrrrrr}
\toprule
Configuration & Skips & Params & FLOPs & Mean & Median & P90\\
\midrule
no-skip, dec$\times$0.5 & no  & 2.02M & 0.92G & 0.881 & 0.721 & 1.565\\
no-skip, dec$\times$1.0 & no  & 2.27M & 1.28G & \textbf{0.882}$\pm$0.018 & 0.719 & 1.570\\
no-skip, dec$\times$2.0 & no  & 2.92M & 2.51G & 0.883 & 0.721 & 1.562\\
\midrule
skip, dec$\times$0.5    & yes & 2.09M & 1.00G & 0.851 & 0.671 & 1.513\\
skip, dec$\times$1.0    & yes & 2.41M & 1.45G & \textbf{0.830}$\pm$0.003 & 0.695 & 1.435\\
skip, dec$\times$2.0    & yes & 3.19M & 2.85G & 0.866 & 0.721 & 1.513\\
\midrule
skip + ASPP, dec$\times$1.0 & yes & 3.28M & --- & \textbf{0.791} & 0.671 & 1.432\\
\bottomrule
\end{tabular}
\end{table}

\paragraph{The headline: capacity alone does nothing; spatial structure does.}
Tripling the no-skip decoder (2.02\,M $\to$ 2.92\,M params, FLOPs $\times$2.7)
leaves the error \emph{flat}: $0.881 \to 0.882 \to 0.883$\,m. Adding skip
connections lowers the error at \emph{every} width. At matched decoder width the
improvement from spatial preservation is:

\begin{center}\small
\begin{tabular}{lrrr}
\toprule
 & no-skip (2.27M) & skip (2.41M) & improvement\\
\midrule
Mean error   & 0.882\,m & 0.830\,m & \textbf{$-$0.052\,m ($-5.9\%$)}\\
P90 (tail)   & 1.570\,m & 1.435\,m & \textbf{$-$0.135\,m ($-8.6\%$)}\\
+ multi-scale (ASPP) mean & 0.882\,m & 0.791\,m & \textbf{$-$0.091\,m ($-10.3\%$)}\\
\bottomrule
\end{tabular}
\end{center}

\noindent The paired per-sample test is decisive (Wilcoxon $p=1.8\times10^{-11}$,
$n=4260$; seed std only $\pm0.003$). The gain is \emph{larger at the tail}
(P90 $-8.6\%$ vs.\ mean $-5.9\%$): skips mainly fix the hard, spatially-ambiguous
samples. A skip model with \emph{fewer} parameters also beats a no-skip model
with \emph{more} (skip@2.41M $=0.830$ vs.\ no-skip@2.92M $=0.883$); at equal
parameter count the skip curve is lower by $0.031$--$0.053$\,m.

\paragraph{Versus the old architecture.} The improvement above is purely
architectural (no-skip vs.\ skip twins). For context against the \emph{old}
pipeline, the full DLoc baseline is a 16.5\,M-parameter ResNet encoder--decoder;
our skip model reaches comparable or better accuracy at \textbf{$\approx$6.8$\times$
fewer parameters} (Table~\ref{tab:oldnew}, medians shown for continuity with the
thesis\footnote{Thesis/paper headline figures are reported as \emph{medians};
Study~A selects and compares by \emph{mean}. The two metrics rank the models
consistently; only the absolute values differ.}).

\begin{table}[h]
\centering
\caption{Old architecture vs.\ new lightweight models (median error).}
\label{tab:oldnew}
\setlength{\tabcolsep}{6pt}\footnotesize
\begin{tabular}{lrrl}
\toprule
Model & Params & Median (m) & Note\\
\midrule
DLoc baseline (old, ResNet enc--dec) & 16.5M & 0.707\footnotemark[2] & big model\\
MobileNetV2, no-skip (new) & 2.27M & 0.719 & $7.3\times$ smaller, no skips\\
MobileNetV2-UNet, skip (new) & 2.41M & \textbf{0.695}\footnotemark[3] & $6.8\times$ smaller, with skips\\
\bottomrule
\end{tabular}
\end{table}
\footnotetext[2]{Our reproduction; multi-seed mean $0.875\pm0.036$\,m. Published
DLoc fig.\,10b complex-space median $\approx0.64$\,m.}
\footnotetext[3]{Study~A best-by-mean checkpoint; the thesis U-Net best-by-median
run reaches $0.640$\,m, tying published DLoc at $6.8\times$ compression.}

%============================================================
\section{Depth of each model}

``Depth'' here is the number of learnable weight layers (Conv/Deconv/Linear)
along the network. \textbf{Crucially, Study~A holds depth constant}: all 13 runs
have the same depth: the ablation varies \emph{width} (channels) and
\emph{skip-connectivity}, never depth. So the accuracy differences in \S1 cannot
be attributed to a deeper network. (Numbers below are reproduced exactly by
\texttt{code/evaluation/model\_depths.py}.)

\begin{table}[h]
\centering
\caption{Model depths and parameter counts.}
\setlength{\tabcolsep}{6pt}\footnotesize
\begin{tabular}{lrrl}
\toprule
Model & Params & Depth & Composition\\
\midrule
DLoc baseline (old) & 16.5M & $\approx$36 & 6 enc + 9 dec ResNet blocks + down/up/IO\\
MobileNetV2, no-skip & 2.27M & $\approx$55 & 51-layer MobileNetV2 encoder + 4-layer decoder\\
MobileNetV2-UNet, skip & 2.41M & $\approx$55 & \emph{same depth} + 3 skip connections\\
MobileNetV2-UNet + ASPP & 3.28M & $\approx$56 & + multi-scale bottleneck module\\
TinyCNN & 47K & $\approx$12 & 3 conv blocks + 2 deconv + 1$\times$1 head\\
Mamba-DLoc & 653K & --- & 3-conv stem + 4 Mamba blocks + 3-conv decoder\\
\bottomrule
\end{tabular}
\end{table}

\noindent The no-skip and skip students are \emph{identical in depth and encoder};
they differ only in whether three high-resolution encoder feature maps are routed
into the decoder. That is the only variable behind the \S1 improvement.

%============================================================
\section{Why skip connections increase accuracy}

The intuition raised (``the gradient to the early layers is very small, so the
skip carries over part of what was learned in the initial layers'') is correct,
and it is actually \emph{two} mechanisms that a skip connection provides at once.

A deep encoder--decoder suffers two things as it goes deep:
\begin{enumerate}
\item \textbf{Backward (gradient flow).} The error gradient must travel back
through the whole bottleneck to reach the early encoder layers; through dozens of
layers it attenuates (the vanishing-gradient effect), so the early layers barely
update. A skip connection opens a \emph{short path}: the loss gradient flows
directly from the decoder into the mid/early encoder features via the
concatenation, bypassing the bottleneck. The early layers then receive a larger,
less-attenuated gradient and keep learning. \emph{This is exactly the
``gradient-is-small $\Rightarrow$ skip helps'' reasoning.}
\item \textbf{Forward (feature reuse).} The encoder downsamples
$161\times361 \to \approx 6\times12$ (a $\sim$27$\times$ spatial reduction), which
destroys fine spatial resolution. The skip copies the \emph{high-resolution}
early feature maps straight into the decoder, so the decoder can reconstruct the
sharp 2D structure --- the precise peak location --- that the bottleneck threw
away. This is literally ``carrying over what was learned in the initial layers,''
but as a forward signal.
\end{enumerate}

For a localization \emph{heatmap} --- where the target is a precise 2D position
--- the forward spatial-detail path (2) is likely the dominant driver, which is
why the gain is largest at the hard tail (P90 $-8.6\%$ vs.\ mean $-5.9\%$ in
\S1). The gradient experiment in \S4 measures the backward path (1) directly.

\paragraph{The testable prediction.} If the gradient argument holds, then
without skips the gradient norm should \emph{collapse} toward the earliest layers,
and skips should \emph{restore/flatten} it. Equivalently, the early-layer
(\texttt{inc}, \texttt{enc1}) gradient norm should be higher \emph{with} skips.

%============================================================
\section{Gradient-descent snapshot and gradient-flow graphs}

The script \texttt{code/evaluation/studyA\_gradients.py} re-trains the two anchor
models (no-skip and skip, decoder width~1.0, seed~42) with the identical setup of
\texttt{train\_studyA.py}, and logs, every epoch, the training/test curves and the
L2 norm of the gradient reaching each layer group.

\begin{figure}[h]
\centering
\figbox{StudyA-figs/fig-training-curves.png}{0.7}
\caption{\textbf{Gradient-descent snapshot.} Test mean error vs.\ epoch, skip vs.\
no-skip. The skip curve descends to a lower error and stays below the no-skip
curve throughout.}
\end{figure}

\begin{figure}[h]
\centering
\figbox{StudyA-figs/fig-grad-profile.png}{0.72}
\caption{\textbf{Gradient vs.\ depth (the key plot).} Gradient L2 norm (log scale)
reaching each layer group, ordered input\,$\to$\,output; early layers highlighted.
The gradient argument predicts the no-skip profile drops toward the early layers
while the skip profile stays flatter --- skips restore the gradient to the
initial layers.}
\end{figure}

\begin{figure}[h]
\centering
\figbox{StudyA-figs/fig-grad-early.png}{0.7}
\caption{\textbf{Early-layer gradient over training.} Gradient norm at the first
conv (\texttt{inc}) vs.\ epoch, skip vs.\ no-skip.}
\end{figure}

\noindent\emph{(Figures populate automatically once the $\sim$1.5--2.2\,h run
finishes and the three PNGs are placed in \texttt{StudyA-figs/}; the one-line
confirmation of the predicted pattern will be added then.)}

%============================================================
\section{Conclusion}

Keeping the spatial relation between adjacent pixels --- via skip connections,
and further via multi-scale (ASPP) --- improves localization accuracy by
\textbf{$\approx$6\% (mean) / $\approx$9\% (P90)} for skips and
\textbf{$\approx$10\%} with multi-scale, \emph{at essentially constant depth and
with no benefit from added parameters} (tripling the plain decoder changed
nothing). The mechanism is a skip connection acting simultaneously as a
gradient highway to the early layers and as a forward high-resolution
feature path; for a 2D localization heatmap the latter dominates, consistent with
the larger tail-error improvement. This directly supports the thesis's choice of
a U-Net student and reframes the Mamba failure as a loss of 2D spatial locality
rather than a lack of capacity.

\end{document}
